\chapter{Описание экспериментов по фаззингу} \label{ch2}

Данная глава содержит результаты запусков фаззера, входные данные, позволившие
выявить ошибки, а также оценку времени, памяти, удобства и эффективности
динамического анализа.

\section{Условия проведения эксперимента} \label{ch2:sec1}

Эксперименты выполнялись в локальном каталоге приложения. Для каждого fuzz target
использовалась команда с ограничением \texttt{-runs=20000}. Фактически оба дефекта
были найдены до достижения установленного лимита запусков.

Условия эксперимента приведены в \taref{tab:fuzz-exp-env}.

\noindent
\begingroup
\centering
\small
\begin{longtable}[c]{|p{5cm}|p{10cm}|}
    \caption{Условия проведения эксперимента по фаззингу}
    \label{tab:fuzz-exp-env}
    \\
    \hline
    Параметр & Значение \\ \hline
    \endfirsthead
    \hline
    Параметр & Значение \\ \hline
    \endhead
    \hline
    Рабочий каталог & \texttt{practice3/palindrome\_app} \\ \hline
    Фаззер & Jazzer.js \\ \hline
    Целевые модули & \texttt{datasetNameParser.js}, \texttt{percentInputDecoder.js} \\ \hline
    Ограничение запусков & \texttt{-runs=20000} \\ \hline
    Сохранение результатов & \texttt{reports/fuzz/dataset-fuzz.txt}, \texttt{reports/fuzz/percent-fuzz.txt} \\ \hline
\end{longtable}
\normalsize
\endgroup

\section{Затраты по времени и памяти} \label{ch2:sec2}

Jazzer.js выводит количество выполненных запусков и показатель \texttt{rss}, то есть
объём используемой памяти процесса. Результаты приведены в \taref{tab:fuzz-costs}.

\noindent
\begingroup
\centering
\small
\begin{longtable}[c]{|p{4.2cm}|p{3cm}|p{3cm}|p{4.5cm}|}
    \caption{Затраты по времени и памяти при фаззинге}
    \label{tab:fuzz-costs}
    \\
    \hline
    Fuzz target & Запусков до ошибки & Память & Итог \\ \hline
    \endfirsthead
    \hline
    Fuzz target & Запусков до ошибки & Память & Итог \\ \hline
    \endhead
    \hline
    \texttt{datasetNameFuzzTarget.js} & 1 & Не выводилась отдельно до аварии & Найден \texttt{TypeError} \\ \hline
    \texttt{percentDecoderFuzzTarget.js} & 134 & 170 МБ & Найден \texttt{URIError} \\ \hline
\end{longtable}
\normalsize
\endgroup

В обоих случаях ошибка была найдена практически сразу. Для \texttt{datasetNameFuzzTarget}
минимальным входом стала пустая строка. Для \texttt{percentDecoderFuzzTarget}
минимальным входом стала строка \texttt{\%}, сохранённая фаззером в crash-файл с
Base64-представлением \texttt{JQ==}.

\section{Итоговая таблица результатов фаззинга} \label{ch2:sec3}

Итоговая информация по найденным дефектам приведена в \taref{tab:fuzz-results}.

\noindent
\begingroup
\centering
\small
\begin{longtable}[c]{|p{3.8cm}|p{3cm}|p{3.5cm}|p{4.7cm}|}
    \caption{Итоговые результаты фаззинг-тестирования}
    \label{tab:fuzz-results}
    \\
    \hline
    Модуль & Подобранные данные & Исключение & Суть ошибки \\ \hline
    \endfirsthead
    \hline
    Модуль & Подобранные данные & Исключение & Суть ошибки \\ \hline
    \endhead
    \hline
    \texttt{datasetNameParser.js} & Пустая строка & \texttt{TypeError} & Нет проверки количества частей имени датасета \\ \hline
    \texttt{percentInputDecoder.js} & \texttt{\%} & \texttt{URIError} & Нет обработки некорректного percent-кодирования \\ \hline
\end{longtable}
\normalsize
\endgroup

\section{Обсуждение результатов} \label{ch2:sec4}

Jazzer.js показал высокую скорость поиска ошибок в выбранных модулях. Оба дефекта
были найдены без ручного подбора входных данных. Детализация сообщения достаточна
для локализации ошибки: вывод содержит тип исключения, стек вызовов, строку файла
и минимальный вход, сохранённый в crash-файл.

Для модуля \texttt{datasetNameParser.js} фаззер показал, что функция не проверяет
структуру входной строки перед обращением к элементу массива. Для модуля
\texttt{percentInputDecoder.js} фаззер показал, что стандартная функция
\texttt{decodeURIComponent} может аварийно завершить обработку некорректного
пользовательского ввода.

\section{Выводы} \label{ch2:conclusion}

Фаззинг оказался удобным и эффективным для поиска ошибок в функциях, работающих
с пользовательским вводом. Для данного приложения он особенно полезен на границах
форматов: имена датасетов, JSON-строки, percent-кодирование и строковые параметры
алгоритма.
